% =====================================================================
%  sections/03_analysis.tex  --  III. ANALYSIS
% =====================================================================
\section{ANALYSIS}
\label{sec:analysis}

% ---------------------------------------------------------------------
\subsection{Scope}
\label{sub:scope}

Project scope outlines the boundaries and defines exactly what goals,
deadlines, and project deliverables to work towards. By clarifying this,
it will assist us in hitting the project goals and objectives without
delay or overwork.

\noindent The specific functional scope implemented for \projname{} is detailed in Table~\ref{tab:scope}.

\begin{table}[H]
\centering
\caption{Functional scope of \protect\projnameplain{}, listing each core feature (account management, dashboard, chatbot, forecasting, report generation, and visual container) with a description of its intended capability.}
\label{tab:scope}
\small
\begin{tabularx}{\linewidth}{L{4.2cm} Y}
\rowcolor{HeaderGray} \textcolor{white}{\textbf{Function}} & \textcolor{white}{\textbf{Description}} \\
Account Management &
  The administrator has full control over all employee accounts within
  the organization, including the ability to edit, delete, activate, and
  deactivate them. \\
\rowcolor{RowGray}
Dashboard &
  Users can see the overview of key data insights through interactive
  charts and KPIs visualization. \\
ChatBot &
  Users can ask the chat AI for insights, including visualizations,
  tables, and readable text. \\
\rowcolor{RowGray}
Sales and Purchase Forecasting &
  Users are able to conduct sales and purchase forecasting to analyze
  trends and predict future patterns. \\
Report Generation &
  Users have the capability to generate quick reports with chart
  visualization and customize them as needed. \\
\rowcolor{RowGray}
Visual Container &
  Users can store visual, markdown, and drag-and-drop visuals in the
  report. \\
\end{tabularx}
\end{table}

% ---------------------------------------------------------------------
\subsection{Risk Management}
\label{sub:risk-management}

Risk management involves identifying and analyzing potential project
risks and minimizing their impact on project progress. To ensure smooth project execution, we have identified the key potential risks that could impact the system, which are detailed in Table~\ref{tab:risks}.

\begin{table}[H]
\centering
\caption{Key project risks identified during planning, including AI hallucination and technical implementation risk, along with a description of each risk's potential impact on the system.}
\label{tab:risks}
\small
\begin{tabularx}{\linewidth}{L{4.2cm} Y}
\rowcolor{HeaderGray} \textcolor{white}{\textbf{Type of Risk}} & \textcolor{white}{\textbf{Description}} \\
AI hallucination &
  The possibility that the AI may generate inaccurate, misleading
  responses that impact data accuracy and project outcome. \\
\rowcolor{RowGray}
Technical Risk &
  The inability to implement advanced technology is due to a lack of
  expertise and/or experience. \\
\end{tabularx}
\end{table}